\medskip

\textbf{Principaux domaines abordés :} 
probabilités.

\medskip

\emph{Cet exercice est un questionnaire à choix multiples.\\
Pour chacune des questions suivantes, une seule des quatre réponses proposées est exacte.\\
Une réponse fausse, une réponse multiple ou l'absence de réponse à une question ne rapporte ni n'enlève de point.\\
Pour répondre, indiquer sur la copie le numéro de la question et la lettre de la réponse choisie. Aucune justification n'est demandée.}

\medskip

On considère un système de communication binaire transmettant des $0$
et des $1$.

Chaque $0$ ou $1$ est appelé bit.

En raison d'interférences, il peut y avoir des erreurs de transmission:

un $0$ peut être reçu comme un $1$ et, de même, un $1$ peut être reçu comme un $0$.

Pour un bit choisi au hasard dans le message, on note les évènements :

\begin{minipage}{0.48\linewidth}

\begin{itemize}
\item[$\bullet~~$] $E_0$ : \og le bit envoyé est un $0$ \fg{} ;
\item[$\bullet~~$] $E_1$ : \og le bit envoyé est un 1 \fg{} ;
\item[$\bullet~~$] $R_0$ : \og le bit reçu est un $0$\fg{} 
\item[$\bullet~~$] $R_1$ : \og le bit reçu est un $1$ \fg.
\end{itemize}
\end{minipage}\hfill
\begin{minipage}{0.48\linewidth}
\begin{center}
\pstree[treemode=R,nodesepA=0pt,nodesepB=2.5pt,treesep = 1cm,levelsep=2.5cm]{\TR{}}
{\pstree{\TR{$E_0$~}\taput{0,4}}
{\TR{$R_0$}\taput{\ldots}
\TR{$R_1$}\tbput{0,01}
}
\pstree{\TR{$E_1$~}\tbput{\ldots}}
{\TR{$R_0$}\taput{0,02}
\TR{$R_1$}\tbput{\ldots}
}}
\end{center}
\end{minipage}

On sait que:

$p\left(E_0\right) = 0,4 \:;\quad p_{E_0}\left(R_1\right) = 0,01 \:;\quad p_{E_1}\left(R_0\right) = 0,02$.

On rappelle que la probabilité conditionnelle de $A$ sachant $B$ est notée $p_B(A)$.

On peut ainsi représenter la situation par l'arbre de probabilités ci-dessus.

\medskip

\begin{enumerate}
\item La probabilité que le bit envoyé soit un $0$ et que le bit reçu soit un $0$ est égale à :

\begin{center}
\begin{tabularx}{\linewidth}{*{4}{X}}
\textbf{a.~~}0,99 &\textbf{b.~~}0,396 &\textbf{c.~~}0,01 &\textbf{d.~~} 0,4
\end{tabularx}
\end{center}

\item La probabilité  $p\left(R_0\right)$ est égale à :
\begin{center}
\begin{tabularx}{\linewidth}{*{4}{X}}
\textbf{a.~~}0,99 &\textbf{b.~~}0,02 &\textbf{c.~~}0,408 &\textbf{d.~~}0,931
\end{tabularx}
\end{center}
\item Une valeur, approchée au millième, de la probabilité $p_{R_1}\left(E_0\right)$ est égale 
\begin{center}
\begin{tabularx}{\linewidth}{*{4}{X}}
\textbf{a.~~}0,004 &\textbf{b.~~}0,001 &\textbf{c.~~}0,007 &\textbf{d.~~}0,010
\end{tabularx}
\end{center}
\item La probabilité de l'évènement \og  il y a une erreur de transmission \fg{} est égale à :
\begin{center}
\begin{tabularx}{\linewidth}{*{4}{X}}
\textbf{a.~~}0,03 &\textbf{b.~~}0,016 &\textbf{c.~~}0,16 &\textbf{d.~~}0,015
\end{tabularx}
\end{center}
\end{enumerate}

Un message de longueur huit bits est appelé un octet.

On admet que la probabilité qu'un octet soit transmis sans erreur est égale à 0,88.

\begin{enumerate}[resume]
\item On transmet successivement 10 octets de façon indépendante.

La probabilité, à $10^{-3}$ près, qu'exactement 7 octets soient transmis sans erreur est égale à :

\begin{center}
\begin{tabularx}{\linewidth}{*{4}{X}}
\textbf{a.~~}0,915 &\textbf{b.~~}0,109 &\textbf{c.~~}0,976 &\textbf{d.~~}0,085
\end{tabularx}
\end{center}
\item On transmet successivement 10 octets de façon indépendante.

La probabilité qu'au moins 1 octet soit transmis sans erreur est égale à :

\begin{center}
\begin{tabularx}{\linewidth}{*{4}{X}}
\textbf{a.~~} $1 - 0,12^{10}$ &\textbf{b.~~}$0,12^{10}$ &\textbf{c.~~}$0,88^{10}$ &\textbf{d.~~} $1- 0,88^{10}$
\end{tabularx}
\end{center}
\item Soit $N$ un entier naturel. On transmet successivement $N$ octets de façon indépendante. 

Soit $N_0$ la plus grande valeur de $N$ pour laquelle la probabilité que les $N$ octets soient tous transmis sans erreur est supérieure ou égale à $0,1$.

On peut affirmer que:

\begin{center}
\begin{tabularx}{\linewidth}{*{4}{X}}
\textbf{a.~~}$N_0 = 17$ &\textbf{b.~~}$N_0 = 18$ &\textbf{c.~~}$N_0 = 19$ &\textbf{d.~~}$N_0 = 20$
\end{tabularx}
\end{center}
\end{enumerate}
